\documentclass[12pt,a4paper]{article}
\usepackage[utf8]{inputenc}
\usepackage[spanish]{babel}
\usepackage{geometry}
\usepackage{fancyhdr}
\usepackage{enumitem}
\usepackage{xcolor}
\usepackage{listings}
\usepackage{graphicx}
\usepackage{hyperref}
\usepackage{tabularx}
\usepackage{array}

\geometry{margin=2cm}
\pagestyle{fancy}
\fancyhf{}
\lhead{Curso de Programación en Python}
\rhead{Semana 5}
\cfoot{\thepage}

% Configuración para código Python
\lstset{
    language=Python,
    basicstyle=\ttfamily\small,
    keywordstyle=\color{blue},
    stringstyle=\color{red},
    commentstyle=\color{green!60!black},
    showstringspaces=false,
    breaklines=true,
    frame=single,
    numbers=left,
    numberstyle=\tiny\color{gray}
}

\title{\textbf{Tarea Semana 5: Calculadora Científica con Tests}\\
\large Pruebas Unitarias y Testing}
\author{Curso de Programación en Python}
\date{Tiempo de entrega: 1 semana}

\begin{document}

\maketitle

\section{Introducción}

¡Bienvenido a la Semana 5! Has aprendido a programar, tomar decisiones, usar bucles, manejar errores y trabajar con librerías. Ahora es momento de aprender una habilidad profesional crucial: \textbf{testing}.

¿Cómo sabes que tu código funciona correctamente? ¿Cómo garantizas que funciona en todos los casos posibles? ¿Cómo te aseguras que un cambio no rompe funcionalidad existente? La respuesta es: \textbf{pruebas unitarias}.

Las pruebas unitarias son código que verifica que tu código funciona correctamente. Los programadores profesionales escriben tests para cada función que crean. En esta tarea, aprenderás a usar \texttt{pytest}, el framework de testing más popular de Python, para crear pruebas exhaustivas de una calculadora científica.

\section{Objetivos de Aprendizaje}

Al completar esta tarea, serás capaz de:

\begin{itemize}[leftmargin=*]
    \item Entender la importancia del testing en desarrollo de software
    \item Instalar y configurar \texttt{pytest}
    \item Escribir funciones de prueba con \texttt{pytest}
    \item Usar la función \texttt{assert} para verificar resultados
    \item Crear pruebas para casos normales, límite y excepciones
    \item Organizar tests en archivos separados
    \item Ejecutar tests y entender reportes de pytest
    \item Aplicar Test-Driven Development (TDD) básico
    \item Calcular cobertura de código con \texttt{pytest-cov}
    \item Identificar y corregir bugs usando tests
\end{itemize}

\section{Enunciado del Proyecto}

\subsection{Descripción General}

Vas a crear una \textbf{Calculadora Científica} completa con operaciones matemáticas avanzadas, Y crear pruebas unitarias exhaustivas para cada función usando \texttt{pytest}.

Este proyecto tiene DOS partes iguales en importancia:
\begin{enumerate}
    \item Implementar las funciones de la calculadora
    \item Escribir tests completos para cada función
\end{enumerate}

\subsection{Filosofía TDD (Test-Driven Development)}

Opcionalmente, puedes aplicar TDD:
\begin{enumerate}
    \item Escribir el test PRIMERO (que falla porque la función no existe)
    \item Implementar la función mínima para que pase
    \item Refactorizar si es necesario
    \item Repetir
\end{enumerate}

\section{Requisitos Funcionales}

\subsection{Parte 1: Funciones a Implementar}

Debes crear un archivo \texttt{calculadora.py} con las siguientes funciones:

\subsubsection{Operaciones Básicas}

\begin{lstlisting}[caption=Operaciones basicas]
def sumar(a, b):
    """Suma dos numeros"""
    return a + b

def restar(a, b):
    """Resta b de a"""
    return a - b

def multiplicar(a, b):
    """Multiplica dos numeros"""
    return a * b

def dividir(a, b):
    """Divide a entre b
    Raises:
        ValueError: Si b es 0
    """
    if b == 0:
        raise ValueError("No se puede dividir entre cero")
    return a / b
\end{lstlisting}

\subsubsection{Operaciones Avanzadas}

\begin{lstlisting}[caption=Operaciones avanzadas]
def potencia(base, exponente):
    """Calcula base elevado a exponente"""
    return base ** exponente

def raiz_cuadrada(numero):
    """Calcula la raiz cuadrada
    Raises:
        ValueError: Si numero es negativo
    """
    if numero < 0:
        raise ValueError("No se puede raiz de numero negativo")
    return numero ** 0.5

def factorial(n):
    """Calcula el factorial de n
    Raises:
        ValueError: Si n es negativo o no es entero
    """
    if not isinstance(n, int):
        raise ValueError("n debe ser un entero")
    if n < 0:
        raise ValueError("n debe ser positivo")
    if n == 0 or n == 1:
        return 1
    resultado = 1
    for i in range(2, n + 1):
        resultado *= i
    return resultado

def es_primo(n):
    """Verifica si n es un numero primo
    Returns:
        bool: True si es primo, False si no
    Raises:
        ValueError: Si n no es entero positivo
    """
    if not isinstance(n, int) or n < 2:
        raise ValueError("n debe ser entero mayor o igual a 2")
    if n == 2:
        return True
    if n % 2 == 0:
        return False
    for i in range(3, int(n**0.5) + 1, 2):
        if n % i == 0:
            return False
    return True
\end{lstlisting}

\subsubsection{Operaciones con Listas}

\begin{lstlisting}[caption=Operaciones con listas]
def promedio(numeros):
    """Calcula el promedio de una lista
    Raises:
        ValueError: Si la lista esta vacia
    """
    if not numeros:
        raise ValueError("La lista no puede estar vacia")
    return sum(numeros) / len(numeros)

def mediana(numeros):
    """Calcula la mediana de una lista
    Raises:
        ValueError: Si la lista esta vacia
    """
    if not numeros:
        raise ValueError("La lista no puede estar vacia")
    ordenados = sorted(numeros)
    n = len(ordenados)
    if n % 2 == 0:
        return (ordenados[n//2 - 1] + ordenados[n//2]) / 2
    return ordenados[n//2]

def maximo(numeros):
    """Retorna el maximo de una lista
    Raises:
        ValueError: Si la lista esta vacia
    """
    if not numeros:
        raise ValueError("La lista no puede estar vacia")
    return max(numeros)

def minimo(numeros):
    """Retorna el minimo de una lista
    Raises:
        ValueError: Si la lista esta vacia
    """
    if not numeros:
        raise ValueError("La lista no puede estar vacia")
    return min(numeros)
\end{lstlisting}

\subsection{Parte 2: Tests a Implementar}

Debes crear un archivo \texttt{test\_calculadora.py} con tests para TODAS las funciones:

\subsubsection{Estructura de un Test}

\begin{lstlisting}[caption=Ejemplo de test basico]
import pytest
from calculadora import sumar, dividir, raiz_cuadrada

# Test simple
def test_sumar_dos_positivos():
    """Verifica suma de dos numeros positivos"""
    resultado = sumar(2, 3)
    assert resultado == 5

# Test con multiples casos
def test_sumar_varios_casos():
    """Verifica suma en varios escenarios"""
    assert sumar(0, 0) == 0
    assert sumar(-1, 1) == 0
    assert sumar(-5, -3) == -8
    assert sumar(100, 200) == 300

# Test de excepciones
def test_dividir_entre_cero_lanza_error():
    """Verifica que dividir entre 0 lanza ValueError"""
    with pytest.raises(ValueError):
        dividir(10, 0)

# Test con mensaje de error especifico
def test_raiz_negativo_mensaje_error():
    """Verifica el mensaje de error especifico"""
    with pytest.raises(ValueError, match="No se puede raiz"):
        raiz_cuadrada(-4)
\end{lstlisting}

\subsubsection{Tests Requeridos para cada Función}

Para CADA función debes crear al menos:

\begin{enumerate}[leftmargin=*]
    \item \textbf{Test de caso normal:} Entrada típica esperada
    \item \textbf{Test de caso límite:} Valores en los bordes (0, 1, -1, máximos, mínimos)
    \item \textbf{Test de excepciones:} Verificar que se lanzan errores apropiados
    \item \textbf{Test de tipos:} Verificar comportamiento con diferentes tipos
\end{enumerate}

\subsubsection{Ejemplo Completo: Tests para \texttt{factorial}}

\begin{lstlisting}[caption=Tests completos para factorial]
import pytest
from calculadora import factorial

# Casos normales
def test_factorial_de_0():
    """Factorial de 0 debe ser 1"""
    assert factorial(0) == 1

def test_factorial_de_1():
    """Factorial de 1 debe ser 1"""
    assert factorial(1) == 1

def test_factorial_de_5():
    """Factorial de 5 debe ser 120"""
    assert factorial(5) == 120

def test_factorial_de_10():
    """Factorial de 10 debe ser 3628800"""
    assert factorial(10) == 3628800

# Casos limite
def test_factorial_numero_grande():
    """Factorial de numero grande"""
    resultado = factorial(20)
    assert resultado == 2432902008176640000

# Casos de error
def test_factorial_negativo_lanza_error():
    """Factorial de negativo debe lanzar ValueError"""
    with pytest.raises(ValueError, match="debe ser positivo"):
        factorial(-5)

def test_factorial_float_lanza_error():
    """Factorial de float debe lanzar ValueError"""
    with pytest.raises(ValueError, match="debe ser un entero"):
        factorial(5.5)

def test_factorial_string_lanza_error():
    """Factorial de string debe lanzar ValueError"""
    with pytest.raises(ValueError):
        factorial("cinco")

# Tests parametrizados (avanzado)
@pytest.mark.parametrize("n, esperado", [
    (0, 1),
    (1, 1),
    (2, 2),
    (3, 6),
    (4, 24),
    (5, 120),
])
def test_factorial_valores_conocidos(n, esperado):
    """Verifica factorial para valores conocidos"""
    assert factorial(n) == esperado
\end{lstlisting}

\newpage

\subsection{Parte 3: Ejecutar Tests}

\subsubsection{Instalación de pytest}

\begin{lstlisting}[language=bash, caption=Instalacion de pytest]
# Instalar pytest
pip install pytest

# Instalar pytest-cov para cobertura
pip install pytest-cov
\end{lstlisting}

\subsubsection{Ejecutar Tests}

\begin{lstlisting}[language=bash, caption=Comandos de pytest]
# Ejecutar todos los tests
pytest

# Ejecutar con mas detalle
pytest -v

# Ejecutar tests de un archivo especifico
pytest test_calculadora.py

# Ejecutar un test especifico
pytest test_calculadora.py::test_sumar_dos_positivos

# Ejecutar con cobertura de codigo
pytest --cov=calculadora --cov-report=html

# Ver que tests pasaron/fallaron
pytest --tb=short
\end{lstlisting}

\subsubsection{Salida Esperada}

\begin{lstlisting}[caption=Ejemplo de salida de pytest]
========================= test session starts =========================
platform win32 -- Python 3.11.5, pytest-7.4.0
collected 45 items

test_calculadora.py::test_sumar_dos_positivos PASSED           [ 2%]
test_calculadora.py::test_sumar_negativos PASSED               [ 4%]
test_calculadora.py::test_restar_basico PASSED                 [ 6%]
test_calculadora.py::test_multiplicar_por_cero PASSED          [ 8%]
test_calculadora.py::test_dividir_normales PASSED              [11%]
test_calculadora.py::test_dividir_entre_cero PASSED            [13%]
test_calculadora.py::test_potencia_basica PASSED               [15%]
test_calculadora.py::test_raiz_cuadrada_4 PASSED               [17%]
test_calculadora.py::test_raiz_negativo_error PASSED           [20%]
test_calculadora.py::test_factorial_5 PASSED                   [22%]
test_calculadora.py::test_factorial_negativo_error PASSED      [24%]
test_calculadora.py::test_es_primo_2 PASSED                    [26%]
test_calculadora.py::test_es_primo_17 PASSED                   [28%]
test_calculadora.py::test_es_primo_4_no_primo PASSED           [31%]
test_calculadora.py::test_promedio_lista PASSED                [33%]
test_calculadora.py::test_promedio_vacio_error PASSED          [35%]
test_calculadora.py::test_mediana_lista_impar PASSED           [37%]
test_calculadora.py::test_mediana_lista_par PASSED             [40%]
... (mas tests) ...

========================== 45 passed in 0.15s =========================

---------- coverage: platform win32, python 3.11.5 -----------
Name              Stmts   Miss  Cover
-------------------------------------
calculadora.py       85      0   100%
-------------------------------------
TOTAL                85      0   100%
\end{lstlisting}

\newpage

\section{Requisitos Técnicos Detallados}

\subsection{Estructura del Proyecto}

\begin{lstlisting}[caption=Estructura de archivos]
calculadora_proyecto/
    calculadora.py           # Funciones de la calculadora
    test_calculadora.py      # Tests con pytest
    requirements.txt         # Dependencias
    README.txt              # Documentacion
    htmlcov/                # Reporte de cobertura (generado)
        index.html
\end{lstlisting}

\subsection{Archivo requirements.txt}

\begin{lstlisting}[caption=requirements.txt]
pytest>=7.4.0
pytest-cov>=4.1.0
\end{lstlisting}

\subsection{Cantidad Mínima de Tests}

Debes crear AL MENOS:

\begin{itemize}[leftmargin=*]
    \item \textbf{40 tests diferentes} (funciones que empiezan con \texttt{test\_})
    \item \textbf{Tests para todas las funciones} implementadas
    \item \textbf{Casos normales:} Al menos 2 por función
    \item \textbf{Casos límite:} Al menos 1 por función
    \item \textbf{Casos de error:} Para funciones que lanzan excepciones
    \item \textbf{Cobertura:} Al menos 90\% de cobertura de código
\end{itemize}

\subsection{Uso de Assert}

\begin{lstlisting}[caption=Diferentes formas de assert]
# Assert simple
assert resultado == 5

# Assert con mensaje personalizado
assert resultado == 5, "La suma deberia dar 5"

# Assert de comparacion
assert resultado > 0
assert resultado >= 10
assert resultado < 100

# Assert de tipo
assert isinstance(resultado, int)
assert isinstance(resultado, float)

# Assert de pertenencia
assert resultado in [1, 2, 3, 4, 5]
assert "error" in str(excepcion)

# Assert de valores aproximados (para floats)
assert abs(resultado - 3.14159) < 0.00001
# O mejor con pytest:
assert resultado == pytest.approx(3.14159)
\end{lstlisting}

\subsection{Uso de pytest.raises}

\begin{lstlisting}[caption=Verificar excepciones]
import pytest

# Verificar que se lanza excepcion
def test_lanza_error():
    with pytest.raises(ValueError):
        funcion_que_falla()

# Verificar mensaje de error
def test_mensaje_error():
    with pytest.raises(ValueError, match="debe ser positivo"):
        factorial(-5)

# Capturar excepcion para mas verificaciones
def test_error_completo():
    with pytest.raises(ValueError) as exc_info:
        dividir(10, 0)
    
    assert "cero" in str(exc_info.value)
\end{lstlisting}

\subsection{Tests Parametrizados (Avanzado - Opcional)}

\begin{lstlisting}[caption=Tests parametrizados con pytest]
import pytest

@pytest.mark.parametrize("a, b, esperado", [
    (2, 3, 5),
    (0, 0, 0),
    (-1, 1, 0),
    (100, 200, 300),
    (-5, -3, -8),
])
def test_sumar_parametrizado(a, b, esperado):
    """Un test que se ejecuta 5 veces con diferentes valores"""
    assert sumar(a, b) == esperado

# Test de es_primo parametrizado
@pytest.mark.parametrize("n", [2, 3, 5, 7, 11, 13, 17, 19, 23])
def test_numeros_primos(n):
    """Verifica que numeros primos son identificados"""
    assert es_primo(n) == True

@pytest.mark.parametrize("n", [4, 6, 8, 9, 10, 12, 15, 16, 18])
def test_numeros_no_primos(n):
    """Verifica que numeros no primos son identificados"""
    assert es_primo(n) == False
\end{lstlisting}

\section{Ejemplo de Test-Driven Development}

\subsection{Paso 1: Escribir el Test Primero}

\begin{lstlisting}[caption=Test primero (falla porque funcion no existe)]
# test_calculadora.py
from calculadora import mcd  # Funcion que aun no existe

def test_mcd_12_y_8():
    """Maximo comun divisor de 12 y 8 debe ser 4"""
    assert mcd(12, 8) == 4

# Al ejecutar: pytest -> ERROR: cannot import 'mcd'
\end{lstlisting}

\subsection{Paso 2: Implementar Función Mínima}

\begin{lstlisting}[caption=Implementacion minima]
# calculadora.py
def mcd(a, b):
    """Calcula el maximo comun divisor"""
    while b:
        a, b = b, a % b
    return a

# Al ejecutar: pytest -> PASSED
\end{lstlisting}

\subsection{Paso 3: Agregar Más Tests}

\begin{lstlisting}[caption=Mas tests para cubrir casos]
def test_mcd_casos_varios():
    assert mcd(12, 8) == 4
    assert mcd(20, 15) == 5
    assert mcd(17, 13) == 1
    assert mcd(100, 50) == 50

def test_mcd_con_cero():
    assert mcd(5, 0) == 5
    assert mcd(0, 5) == 5
\end{lstlisting}

\section{Consejos y Recomendaciones}

\begin{itemize}[leftmargin=*]
    \item \textbf{Un test, un concepto:} Cada test debe verificar UNA cosa específica
    
    \item \textbf{Nombres descriptivos:} \texttt{test\_factorial\_de\_numero\_negativo\_lanza\_error}
    
    \item \textbf{Documenta los tests:} Usa docstrings para explicar qué verifica cada test
    
    \item \textbf{Ejecuta tests frecuentemente:} Después de cada cambio
    
    \item \textbf{Tests independientes:} Cada test debe poder ejecutarse solo
    
    \item \textbf{Casos de error importan:} Verifica que los errores se manejan correctamente
    
    \item \textbf{Cobertura no es todo:} 100\% cobertura no significa código perfecto
    
    \item \textbf{Red-Green-Refactor:} Test falla (red) → Código pasa (green) → Mejorar (refactor)
    
    \item \textbf{Automatiza:} Los tests deben ejecutarse automáticamente
\end{itemize}

\section{Entregables}

\begin{enumerate}[leftmargin=*]
    \item \texttt{calculadora.py} - Todas las funciones implementadas
    \item \texttt{test\_calculadora.py} - Al menos 40 tests
    \item \texttt{requirements.txt} - Dependencias del proyecto
    \item \texttt{README.txt} - Instrucciones de instalación y ejecución
    \item \texttt{reporte\_cobertura.txt} - Captura del reporte de cobertura
    \item TODOS los tests deben pasar (0 fallos)
    \item Cobertura mínima: 90\%
\end{enumerate}

\section{Defensa del Código (60\% de la nota final)}

\subsection{¿Qué es la defensa del código?}

La defensa del código es una reunión virtual donde presentarás y explicarás tu trabajo. Esta defensa es \textbf{obligatoria} y representa el \textbf{60\% de la nota final} de esta tarea.

\subsection{¿Cómo funciona?}

Durante la defensa deberás:

\begin{enumerate}[leftmargin=*]
    \item \textbf{Compartir tu pantalla} mostrando código y terminal.
    
    \item \textbf{Ejecutar los tests en vivo:}
    \begin{itemize}
        \item \texttt{pytest -v} para ver todos los tests
        \item \texttt{pytest --cov=calculadora} para ver cobertura
        \item Demostrar que TODOS los tests pasan
        \item Abrir reporte HTML de cobertura
    \end{itemize}
    
    \item \textbf{Explicar tus tests:}
    \begin{itemize}
        \item Por qué creaste cada tipo de test
        \item Qué verifica cada assert
        \item Cómo elegiste los casos de prueba
        \item Por qué ciertos valores son casos límite
    \end{itemize}
    
    \item \textbf{Demostrar TDD (si aplicaste):}
    \begin{itemize}
        \item Escribir un nuevo test que falle
        \item Implementar la función para que pase
        \item Mostrar el proceso red-green-refactor
    \end{itemize}
    
    \item \textbf{Responder preguntas como:}
    \begin{itemize}
        \item ``¿Por qué es importante hacer testing?''
        \item ``¿Qué diferencia hay entre un caso normal y un caso límite?''
        \item ``¿Cómo funciona pytest.raises?''
        \item ``¿Qué significa 90\% de cobertura?''
        \item ``¿Rompe un test intencionalmente, qué pasa?''
        \item ``¿Cómo verificas excepciones con mensajes específicos?''
        \item ``¿Qué es TDD y cuáles son sus ventajas?''
    \end{itemize}
    
    \item \textbf{Demostrar comprensión de:}
    \begin{itemize}
        \item Propósito de las pruebas unitarias
        \item Estructura de un test (Arrange-Act-Assert)
        \item Tipos de casos de prueba
        \item Interpretación de reportes de pytest
        \item Concepto de cobertura de código
    \end{itemize}
\end{enumerate}

\subsection{Preguntas Típicas de la Defensa}

Prepárate para responder:

\begin{itemize}[leftmargin=*]
    \item ¿Por qué hacer testing? ¿No basta con probar manualmente?
    \item ¿Qué es un test unitario?
    \item ¿Qué hace \texttt{assert}?
    \item ¿Cómo se usa \texttt{pytest.raises}?
    \item ¿Qué es un caso límite? Dame ejemplos
    \item ¿Qué significa que un test pase?
    \item ¿Qué es cobertura de código?
    \item ¿100\% cobertura significa código sin bugs?
    \item ¿Qué es TDD?
    \item ¿Cómo instalaste pytest?
\end{itemize}

\subsection{Consejos para una Buena Defensa}

\begin{itemize}[leftmargin=*]
    \item \textbf{Ejecuta antes:} Verifica que todos los tests pasan antes de la defensa
    \item \textbf{Conoce tus tests:} Entiende cada test que escribiste
    \item \textbf{Practica romper tests:} Cambia código y observa qué tests fallan
    \item \textbf{Entiende pytest:} Lee la documentación básica
    \item \textbf{Cobertura visual:} Abre el HTML de cobertura y navégalo
\end{itemize}

\subsection{Distribución de la Nota}

\begin{itemize}
    \item \textbf{Código funcional (40\%):} Funciones y tests implementados correctamente
    \item \textbf{Defensa del código (60\%):} Tu capacidad de explicar testing
\end{itemize}

\textbf{Importante:} Si no asistes a la defensa o no puedes explicar tus tests, la nota máxima será 40/100.

\section{Desafíos Opcionales (Puntos Extra)}

\begin{enumerate}[leftmargin=*]
    \item \textbf{Tests parametrizados (+5 pts):}
    \begin{itemize}
        \item Usar \texttt{@pytest.mark.parametrize} en al menos 5 tests
        \item Demostrar comprensión de parametrización
    \end{itemize}
    
    \item \textbf{Fixtures de pytest (+6 pts):}
    \begin{itemize}
        \item Crear fixtures con \texttt{@pytest.fixture}
        \item Usar fixtures para datos de prueba reutilizables
        \item Ejemplo: fixture con lista de números para tests de promedio/mediana
    \end{itemize}
    
    \item \textbf{Más funciones matemáticas (+7 pts):}
    \begin{itemize}
        \item Implementar: MCD, MCM, fibonacci, es\_perfecto
        \item Crear tests completos para cada una (10+ tests por función)
    \end{itemize}
    
    \item \textbf{Cobertura 100\% (+4 pts):}
    \begin{itemize}
        \item Lograr exactamente 100\% de cobertura de código
        \item Incluir captura de pantalla del reporte
    \end{itemize}
    
    \item \textbf{Pytest con markers (+5 pts):}
    \begin{itemize}
        \item Usar \texttt{@pytest.mark.slow} para tests lentos
        \item Usar \texttt{@pytest.mark.skip} con razón
        \item Ejecutar solo ciertos markers
    \end{itemize}
    
    \item \textbf{Tests de integración (+6 pts):}
    \begin{itemize}
        \item Crear función que usa múltiples operaciones
        \item Tests que verifican interacción entre funciones
        \item Ejemplo: calculadora de ecuación cuadrática
    \end{itemize}
\end{enumerate}

\textit{Nota: Los puntos extra solo se otorgan si la implementación es correcta Y puedes explicarla en la defensa.}

\newpage

\section{Rúbrica de Evaluación}

\subsection{Distribución General}

\begin{itemize}
    \item \textbf{Código Funcional:} 40\% (40 puntos)
    \item \textbf{Defensa del Código:} 60\% (60 puntos)
    \item \textbf{Total:} 100 puntos
    \item \textbf{Puntos Extra:} Hasta +33 puntos adicionales
\end{itemize}

\subsection{Parte 1: Código Funcional (40 puntos)}

\subsubsection{1.1 Implementación de Funciones (15 puntos)}

\begin{itemize}[leftmargin=*]
    \item \textbf{Operaciones básicas (4 pts):}
    \begin{itemize}
        \item 4 pts: Suma, resta, multiplicación, división correctas
        \item 3 pts: 3 de 4 correctas
        \item 2 pts: 2 de 4 correctas
        \item 0 pts: Menos de 2 correctas
    \end{itemize}
    
    \item \textbf{Operaciones avanzadas (6 pts):}
    \begin{itemize}
        \item 6 pts: Potencia, raíz, factorial, es\_primo correctas
        \item 5 pts: 3 de 4 correctas
        \item 3 pts: 2 de 4 correctas
        \item 0 pts: Menos de 2 correctas
    \end{itemize}
    
    \item \textbf{Operaciones con listas (5 pts):}
    \begin{itemize}
        \item 5 pts: Promedio, mediana, máximo, mínimo correctas
        \item 4 pts: 3 de 4 correctas
        \item 2 pts: 2 de 4 correctas
        \item 0 pts: Menos de 2 correctas
    \end{itemize}
\end{itemize}

\subsubsection{1.2 Calidad de Tests (25 puntos)}

\begin{itemize}[leftmargin=*]
    \item \textbf{Cantidad de tests (5 pts):}
    \begin{itemize}
        \item 5 pts: 40+ tests diferentes
        \item 4 pts: 30-39 tests
        \item 3 pts: 20-29 tests
        \item 1 pt: 10-19 tests
        \item 0 pts: Menos de 10 tests
    \end{itemize}
    
    \item \textbf{Cobertura de casos (8 pts):}
    \begin{itemize}
        \item 8 pts: Tests cubren casos normales, límite y errores para todas las funciones
        \item 6 pts: Buena cobertura pero falta en algunas funciones
        \item 4 pts: Cobertura básica
        \item 0 pts: Cobertura insuficiente
    \end{itemize}
    
    \item \textbf{Uso correcto de assert (4 pts):}
    \begin{itemize}
        \item 4 pts: Todos los assert son apropiados y verifican correctamente
        \item 3 pts: Mayoría correctos
        \item 1 pt: Uso básico
        \item 0 pts: No usa assert correctamente
    \end{itemize}
    
    \item \textbf{Tests de excepciones (4 pts):}
    \begin{itemize}
        \item 4 pts: Usa pytest.raises correctamente para verificar errores
        \item 3 pts: Usa en algunos casos
        \item 1 pt: Uso incorrecto
        \item 0 pts: No verifica excepciones
    \end{itemize}
    
    \item \textbf{Cobertura de código (4 pts):}
    \begin{itemize}
        \item 4 pts: 90\%+ de cobertura
        \item 3 pts: 80-89\% de cobertura
        \item 2 pts: 70-79\% de cobertura
        \item 1 pt: 60-69\% de cobertura
        \item 0 pts: Menos de 60\%
    \end{itemize}
\end{itemize}

\subsection{Parte 2: Defensa del Código (60 puntos)}

\subsubsection{2.1 Comprensión de Testing (30 puntos)}

\begin{itemize}[leftmargin=*]
    \item \textbf{Concepto de testing (8 pts):}
    \begin{itemize}
        \item 8 pts: Explica perfectamente por qué hacer testing y sus beneficios
        \item 6 pts: Explicación correcta con pequeñas dudas
        \item 3 pts: Comprensión básica
        \item 0 pts: No entiende el propósito del testing
    \end{itemize}
    
    \item \textbf{Tipos de casos de prueba (8 pts):}
    \begin{itemize}
        \item 8 pts: Explica claramente casos normales, límite y excepciones con ejemplos
        \item 6 pts: Explica correctamente con alguna confusión
        \item 3 pts: Explicación superficial
        \item 0 pts: No diferencia tipos de casos
    \end{itemize}
    
    \item \textbf{Uso de pytest (8 pts):}
    \begin{itemize}
        \item 8 pts: Explica assert, pytest.raises, ejecución de tests
        \item 6 pts: Explica correctamente lo básico
        \item 3 pts: Conocimiento limitado
        \item 0 pts: No entiende pytest
    \end{itemize}
    
    \item \textbf{Cobertura de código (6 pts):}
    \begin{itemize}
        \item 6 pts: Explica qué es, cómo medirla, qué significa
        \item 4 pts: Comprensión básica
        \item 2 pts: Explicación vaga
        \item 0 pts: No entiende cobertura
    \end{itemize}
\end{itemize}

\subsubsection{2.2 Demostración Práctica (20 puntos)}

\begin{itemize}[leftmargin=*]
    \item \textbf{Ejecución de tests (8 pts):}
    \begin{itemize}
        \item 8 pts: Ejecuta tests con diferentes opciones, todos pasan
        \item 6 pts: Ejecuta correctamente pero algún test falla
        \item 3 pts: Dificultades para ejecutar
        \item 0 pts: No puede ejecutar tests
    \end{itemize}
    
    \item \textbf{Reporte de cobertura (4 pts):}
    \begin{itemize}
        \item 4 pts: Genera y explica reporte de cobertura HTML
        \item 3 pts: Genera reporte con ayuda
        \item 1 pt: Dificultades
        \item 0 pts: No puede generar reporte
    \end{itemize}
    
    \item \textbf{Respuesta a preguntas técnicas (8 pts):}
    \begin{itemize}
        \item 8 pts: Responde correctamente todas las preguntas sobre testing
        \item 6 pts: Responde la mayoría correctamente
        \item 3 pts: Respuestas parciales
        \item 0 pts: No puede responder preguntas básicas
    \end{itemize}
\end{itemize}

\subsubsection{2.3 Calidad y Organización (10 puntos)}

\begin{itemize}[leftmargin=*]
    \item \textbf{Nombres de tests descriptivos (3 pts):}
    \begin{itemize}
        \item 3 pts: Todos los tests tienen nombres claros y descriptivos
        \item 2 pts: Mayoría con buenos nombres
        \item 1 pt: Nombres poco claros
        \item 0 pts: Nombres confusos
    \end{itemize}
    
    \item \textbf{Documentación de tests (4 pts):}
    \begin{itemize}
        \item 4 pts: Todos los tests con docstrings explicativos
        \item 3 pts: Mayoría documentados
        \item 1 pt: Documentación mínima
        \item 0 pts: Sin documentación
    \end{itemize}
    
    \item \textbf{Organización del código (3 pts):}
    \begin{itemize}
        \item 3 pts: Tests bien organizados, agrupados lógicamente
        \item 2 pts: Organización aceptable
        \item 1 pt: Poco organizado
        \item 0 pts: Desorganizado
    \end{itemize}
\end{itemize}

\subsection{Puntos Extra (Hasta +33 puntos)}

\begin{itemize}[leftmargin=*]
    \item \textbf{Tests parametrizados (+5 pts):} Uso correcto de \texttt{@pytest.mark.parametrize}
    \item \textbf{Fixtures (+6 pts):} Creación y uso de fixtures para datos de prueba
    \item \textbf{Más funciones matemáticas (+7 pts):} 4+ funciones adicionales con tests
    \item \textbf{Cobertura 100\% (+4 pts):} Cobertura completa del código
    \item \textbf{Pytest markers (+5 pts):} Uso de markers (slow, skip, etc.)
    \item \textbf{Tests de integración (+6 pts):} Tests que verifican múltiples funciones
\end{itemize}

\textbf{Nota:} Los puntos extra solo se otorgan si la implementación es correcta Y el estudiante puede explicarla completamente durante la defensa.

\subsection{Penalizaciones}

\begin{itemize}[leftmargin=*]
    \item \textbf{-100\% de la defensa:} No asiste sin justificación (nota máxima = 40/100)
    \item \textbf{-50\% de la defensa:} Evidencia de no entender los tests que escribió
    \item \textbf{-10 pts:} Algún test falla al ejecutar
    \item \textbf{-5 pts:} Entrega tardía de 1 día
    \item \textbf{-10 pts:} Entrega tardía de 2-3 días
    \item \textbf{-20 pts:} Entrega tardía de 4+ días
    \item \textbf{-5 pts:} No incluye requirements.txt
    \item \textbf{-5 pts:} Menos de 40 tests
    \item \textbf{-3 pts:} No incluye reporte de cobertura
\end{itemize}

\subsection{Nota Mínima Aprobatoria}

Para aprobar esta tarea necesitas obtener al menos \textbf{60 puntos de 100}.

\subsection{Escala de Calificación}

\begin{itemize}
    \item \textbf{90-100+ pts:} Excelente - Dominio completo de testing
    \item \textbf{80-89 pts:} Muy Bueno - Comprensión sólida de pruebas unitarias
    \item \textbf{70-79 pts:} Bueno - Comprensión aceptable, necesita práctica
    \item \textbf{60-69 pts:} Suficiente - Comprensión básica, necesita refuerzo
    \item \textbf{0-59 pts:} Insuficiente - No aprobado
\end{itemize}

\section{Recursos de Apoyo}

\begin{itemize}[leftmargin=*]
    \item \textbf{Documentación pytest:} \url{https://docs.pytest.org/}
    \item \textbf{Tutorial pytest:} \url{https://docs.pytest.org/en/stable/getting-started.html}
    \item \textbf{Pytest-cov:} \url{https://pytest-cov.readthedocs.io/}
    \item \textbf{Testing en Python:} \url{https://realpython.com/pytest-python-testing/}
    \item \textbf{TDD con Python:} \url{https://realpython.com/python-testing/}
\end{itemize}

\section{Preguntas Frecuentes}

\textbf{P: ¿Por qué hacer testing si puedo probar manualmente?}\\
R: Los tests automáticos son más rápidos, consistentes, y detectan regresiones cuando cambias código.

\textbf{P: ¿Cómo instalo pytest?}\\
R: \texttt{pip install pytest}

\textbf{P: ¿Cómo ejecuto un solo test?}\\
R: \texttt{pytest test\_calculadora.py::test\_nombre\_especifico}

\textbf{P: ¿Qué significa que un test pasa?}\\
R: Que el assert se cumplió, el resultado fue el esperado.

\textbf{P: ¿Cómo veo la cobertura?}\\
R: \texttt{pytest --cov=calculadora --cov-report=html}

\textbf{P: ¿100\% cobertura significa sin bugs?}\\
R: No. Significa que todas las líneas se ejecutaron, pero puede haber casos no probados.

\textbf{P: ¿Qué es un caso límite?}\\
R: Valores en los bordes del rango válido: 0, 1, -1, máximos, mínimos, listas vacías, etc.

\textbf{P: ¿Cuántos tests son suficientes?}\\
R: No hay número mágico. Debes cubrir casos normales, límite y errores para cada función.

\vspace{1cm}

\begin{center}
\textbf{¡Éxito con tus pruebas unitarias!}\\
\textit{Los tests son la red de seguridad del programador.\\
Aprende a escribirlos bien y programarás con confianza.}
\end{center}

\end{document}
